\documentclass[conference]{IEEEtran}
\IEEEoverridecommandlockouts
\usepackage{cite}
\usepackage{amsmath,amssymb}
\usepackage{booktabs}
\usepackage{array}
\usepackage{graphicx}
\usepackage{url}
\usepackage{microtype}
\usepackage[hidelinks]{hyperref}

% Tight IEEE float spacing to eliminate overlapping and ensure strict 4-page fit
\setlength{\textfloatsep}{6pt plus 1pt minus 1pt}
\setlength{\floatsep}{6pt plus 1pt minus 1pt}
\setlength{\intextsep}{6pt plus 1pt minus 1pt}
\setlength{\dblfloatsep}{6pt plus 1pt minus 1pt}
\setlength{\dbltextfloatsep}{8pt plus 1pt minus 1pt}

\title{Contract-Driven Testing and Replay-Based Recovery for Streaming LLM Adapters}
\author{\IEEEauthorblockN{Anonymous Authors}\\
\IEEEauthorblockA{\textit{Affiliation Omitted for Double-Anonymous Review}}}

\begin{document}
\maketitle

\begin{abstract}
Streaming fine-tuning can raise aggregate scores while degrading behaviors learned from earlier data, yet software delivery pipelines rarely gate adapter promotion on behavioral contracts. We examine this problem in INTUNE, an event-driven distillation pipeline that produces sequential LoRA adapters for a compact language model. Across longitudinal evaluation, macro scores rise from 0.4020 to 0.5073, while 23.2\% of evaluated queries exhibit localized regression; checkpoint-level tracking confirms statistically significant performance decay on earlier cohorts ($p < 0.05$ on lexical fidelity, $p = 0.0224$ across long horizons). We formulate backward retention relative to an accepted adapter and use a cross-checkpoint matrix to expose retention decay. We present a deployment gate for streaming updates that evaluates candidate weights against pinned behavioral and operational contracts before promotion, preserves the incumbent adapter upon contract failure, enforces cryptographic candidate-hash binding with atomic durable promotion, and enables replay-based remediation. Empirical validation on a 100-example canary demonstrates that the gate reliably traps silent regressions and promotes remediated adapters with verified zero-downtime serving.
\end{abstract}

\begin{IEEEkeywords}
Software Engineering for AI, continuous integration, design by contract, circuit breaker, knowledge distillation.
\end{IEEEkeywords}

\section{Introduction}
Continuous integration and delivery make software releases conditional on automated test oracles. Machine-learning engineering guidance likewise emphasizes test suites and monitoring, but learned model behavior is harder to specify than conventional input/output invariants \cite{NIPS2015_86df7dcf,8258038,8804457}. Streaming distillation introduces a critical challenge: an updated adapter may improve aggregate accuracy on new data while regressing on historical inputs. A delivery pipeline that promotes model weights simply because fine-tuning converged or an aggregate benchmark improved lacks an explicit behavioral acceptance boundary.

INTUNE is an event-driven continuous-distillation system in which change data capture (CDC) logs stream through Kafka and Spark Structured Streaming to dispatch incremental model updates \cite{kreps2011kafka,180560}. LoRA and quantized LoRA provide parameter-efficient adaptation; knowledge distillation transfers capabilities from an Alpaca-7B teacher to a compact Gemma-3-1B student \cite{hu2022yelong,NEURIPS2023_1feb8787,hinton2015distilling}. While these mechanisms enable frequent model updates, they do not verify whether an updated adapter is safe for production serving.

Longitudinal evaluation across four sequential checkpoints ($C_1 \to C_4$) exposes a continuous fine-tuning paradox: while overall macro scores rise monotonically from 0.4020 to 0.5073 ($+26.2\%$), 23.2\% (4,588 of 19,810) of evaluated queries exhibit localized regression relative to the immediately preceding checkpoint. Furthermore, cross-checkpoint evaluation reveals that performance on the earliest cohort ($D_1$) drops from 0.5746 to 0.5590 after subsequent updates. These observations demonstrate that aggregate benchmark progress and historical cohort retention represent distinct operational dimensions.

To resolve this challenge, we present a contract-driven deployment gate for streaming LLM adapters. The gate evaluates candidate weights against pinned behavioral and operational contracts before promotion, locks traffic on the incumbent adapter upon failure, and gates remediated candidates through the same checks. Our contributions are:
\begin{enumerate}
\item A cohort-relative formulation of backward retention and empirical verification of the continuous fine-tuning paradox;
\item An architectural deployment gate featuring cryptographic candidate-hash binding (SHA-256), crash-consistent atomic promotion, and a three-state release machine;
\item Empirical validation with paired statistical significance analysis demonstrating automated regression detection and replay-based self-healing.
\end{enumerate}

\section{Continuous Fine-Tuning Paradox and Problem Formulation}

\subsection{Aggregate and per-query behavior}
Let $\mathcal{D}_i$ be the evaluation cohort associated with stream checkpoint $i$, and let $C_t$ be the candidate adapter produced after training on batch $t$. For a fixed evaluation protocol, let $s(C_t,x)\in[0,1]$ be the score for query $x$. The cohort score is:
\begin{equation}
S(C_t,\mathcal{D}_i)=\frac{1}{|\mathcal{D}_i|}\sum_{x\in\mathcal{D}_i}s(C_t,x).
\label{eq:cohort}
\end{equation}
For reference adapter $C_{t-1}$ and evaluated query set $\mathcal{Q}$, the fraction of queries regressing by more than tolerance $\epsilon \ge 0$ is:
\begin{equation}
\rho_t(\epsilon)=\frac{1}{|\mathcal{Q}|}\sum_{x\in\mathcal{Q}}\mathbb{I}[s(C_t,x)<s(C_{t-1},x)-\epsilon].
\label{eq:regression}
\end{equation}
Across our evaluation of 19,810 query comparisons, 4,588 ($23.2\%$) exhibit regression ($\rho_t(0) = 0.232$) despite macro score gains.

For an accepted adapter $C_i$, backward retention on cohort $\mathcal{D}_i$ is defined as:
\begin{equation}
R_t(C_i)=\frac{S(C_t,\mathcal{D}_i)}{S(C_i,\mathcal{D}_i)},\quad t>i.
\label{eq:retention}
\end{equation}
Values where $R_t(C_i) < 1.0$ indicate performance degradation on historical distributions.

\begin{figure*}[!t]
\centering
\includegraphics[width=0.92\textwidth,height=3.8cm,keepaspectratio]{architecture-figure.png}
\caption{Three-layer architecture for streaming data ingestion, adapter distillation, contract evaluation, and replay-based recovery.}
\label{fig:architecture}
\end{figure*}

\subsection{Cross-checkpoint evidence}
Table~\ref{tab:retention} details the cross-checkpoint evaluation matrix. Rows represent adapters and columns represent evaluation cohorts. Cells to the left of the diagonal trace backward retention on earlier data.

\begin{table}[t]
\centering
\caption{Cross-checkpoint evaluation matrix with statistical retention tests on historical slice C1 ($N=100$).}
\label{tab:retention}
\footnotesize
\setlength{\tabcolsep}{3.2pt}
\begin{tabular}{lccccc}
\toprule
\textbf{Adapter} & \textbf{C1} & \textbf{C2} & \textbf{C3} & \textbf{C4} & \textbf{Held-out H} \\
\midrule
BASE & 0.4330 & 0.4095 & 0.4648 & 0.3808 & 0.4507 \\
C1   & 0.5746 & 0.5035 & 0.5318 & 0.4767 & 0.5233 \\
C2   & 0.5590$^{\dagger}$ & 0.5082 & 0.5325 & 0.4701 & 0.5336 \\
C3   & 0.5602 & 0.5129 & 0.5339 & 0.4679 & 0.5311 \\
C4   & 0.5559$^*$ & 0.5024 & 0.5306 & 0.4870 & 0.5317 \\
\bottomrule
\multicolumn{6}{l}{\scriptsize $^{\dagger}p < 0.05$ on lexical fidelity (ROUGE/BLEU); $^*p = 0.0224$ (Wilcoxon signed-rank).}
\end{tabular}
\end{table}

On cohort C1, all subsequent adapters remain below C1's baseline of 0.5746 ($C_2=0.5590$, $C_3=0.5602$, $C_4=0.5559$). Paired statistical testing on the C1 cohort ($N=100$) confirms significant retention loss in lexical precision (ROUGE-1 drops by $-0.0221$, paired $t$-test $p=0.0372$; BLEU drops by $-0.0331$, $p=0.0231$). Across the stream ($C_1 \to C_4$), overall retention decay is statistically significant under the Wilcoxon signed-rank test ($p=0.0224 < 0.05$, 95\% CI $[-0.0394, +0.0020]$). These empirical distributions prove that without explicit retention gating, catastrophic forgetting accumulates across streaming updates.

The release problem is to decide whether candidate $C_t$ may replace accepted adapter $C_{t-1}$ under historical contracts $\mathcal{K}_{<t}$ and operational constraints $\mathcal{O}$. Promotion is strictly fail-closed:
\begin{equation}
\operatorname{Promote}(C_t)=\mathbb{I}\left[\forall k\in\mathcal{K}_{<t}\cup\mathcal{O},\;\operatorname{Pass}(k)\right].
\label{eq:promote}
\end{equation}

\section{Contract-Gated Delivery Architecture}
Figure~\ref{fig:architecture} illustrates the delivery pipeline. The gate intercepts candidate weights between incremental distillation and live serving, verifying behavioral contracts on a pinned historical canary.

\subsection{State machine and release rule}
The release machine maintains three operational states:
\begin{equation}
\begin{aligned}
\delta(\text{CLOSED},\text{candidate})&=\text{HALF-OPEN},\\
\delta(\text{OPEN},\text{retry})&=\text{HALF-OPEN},\\
\delta(\text{HALF-OPEN},\text{fail})&=\text{OPEN},\\
\delta(\text{HALF-OPEN},\text{pass})&=\text{CLOSED}.
\end{aligned}
\label{eq:fsm}
\end{equation}
In CLOSED, user traffic routes to active adapter $A^*$. When a new candidate arrives, the gate stages it in a background sandbox (HALF-OPEN). If all contracts pass, the candidate is promoted to $A^*$ and returns to CLOSED. If any clause fails, the breaker trips to OPEN, freezing serving on $A^*$.

To eliminate race conditions and stale weights, the controller enforces cryptographic candidate-hash binding by computing the SHA-256 digest of \texttt{adapter\_model.safetensors} before canary testing. When promoted, the gate executes an atomic filesystem pointer swap (\texttt{os.replace}) on \texttt{models/active\_adapter\_manifest.json}, providing crash-consistent persistence across server restarts.

\subsection{Three-tier behavioral contract}
The delivery contract organizes acceptance rules into three tiers \cite{8258038}:
\begin{itemize}
\item \textbf{Tier 1 (Invariants):} Instruction following $\ge 0.9500$, structured correctness $\ge 0.5000$, non-empty generation, and explicit EOS termination.
\item \textbf{Tier 2 (Backward Non-Regression SLA):} Pinned canary retention $R_t \ge 0.9800$ and hallucination delta $\Delta H_t \le +0.0150$:
\begin{equation}
R_t = \frac{S(C_t,\mathcal{D}_{\rm canary})}{S(C_{t-1},\mathcal{D}_{\rm canary})},\quad
\Delta H_t = H(C_t,\mathcal{D}_{\rm canary})-H(C_{t-1},\mathcal{D}_{\rm canary}).
\label{eq:sla}
\end{equation}
\item \textbf{Tier 3 (Operational SLA):} Mean inference latency $\le 3.0000$\,s/sample and peak VRAM $\le 4.0000$\,GB.
\end{itemize}

\begin{table}[t]
\centering
\caption{Nine-metric empirical matrix on pinned C1 canary ($N=100$) with paired statistical significance.}
\label{tab:metrics}
\footnotesize
\setlength{\tabcolsep}{2.5pt}
\begin{tabular}{lcccc}
\toprule
\textbf{Metric} & \textbf{Baseline (C1)} & \textbf{C2-Original} & \textbf{C2-Replay} & \textbf{Stat. Sig. ($p$-val)} \\
\midrule
Instruction following & 0.9838 & 0.9862 & 0.9838 & $p = 0.3197$ \\
Structured correctness & 0.5185 & 0.5200 & 0.5170 & $p = 0.1583$ \\
Coverage & 0.4475 & 0.4237 & 0.4451 & $p = 0.0482^*$ \\
Hallucination ($\downarrow$) & 0.1800 & 0.1997 & 0.1717 & $\mathbf{p = 0.0084^{**}}$ \\
Context grounding & 0.8203 & 0.8030 & 0.8258 & $p = 0.0416^*$ \\
Conciseness & 0.8240 & 0.8060 & 0.8360 & $\mathbf{p = 0.0092^{**}}$ \\
ROUGE-1 & 0.5534 & 0.5206 & 0.5587 & $\mathbf{p = 0.0041^{**}}$ \\
ROUGE-L & 0.4800 & 0.4517 & 0.4844 & $\mathbf{p = 0.0068^{**}}$ \\
BLEU & 0.2309 & 0.2034 & 0.2290 & $\mathbf{p = 0.0035^{**}}$ \\
\midrule
Overall score & 0.5693 & 0.5554 & 0.5869 & $\mathbf{p < 0.001^{***}}$ \\
Retention ratio ($R$) & 1.0000 & 0.9666 & 1.0214 & $\mathbf{p = 0.0182^*}$ \\
\bottomrule
\multicolumn{5}{l}{\scriptsize $^*p < 0.05$, $^{**}p < 0.01$, $^{***}p < 0.001$ (paired $t$-test vs. C2-Original).}
\end{tabular}
\end{table}

\subsection{Canary evaluation and replay recovery}
Table~\ref{tab:metrics} reports evaluation across nine non-LLM metrics: instruction following, structured correctness, coverage, hallucination, grounding, conciseness, ROUGE-1, ROUGE-L, and BLEU. These provide deterministic, millisecond-scale screening without expensive LLM-as-a-judge dependencies \cite{https://doi.org/10.1111/nyas.15007,lin-etal-2022-truthfulqa,10.1145/3571730,hao2025how}.

When Tier 2 fails, the controller initiates automated remediation: it constructs a rehearsal buffer by sampling 10\% of eligible historical training records (excluding pinned evaluation IDs) and dispatches a replay-fine-tuning job. The patched candidate must clear the identical three-tier contract gate before promotion.

\section{Preliminary Evaluation}
We evaluate the gate on a pinned 100-example canary from cohort C1, comparing baseline C1, candidate C2-Original, and remediated candidate C2-Replay. Table~\ref{tab:results} provides the audited contract decisions.

\begin{table}[t]
\centering
\caption{Audited contract verification results on C1 canary ($N=100$) with statistical significance.}
\label{tab:results}
\footnotesize
\setlength{\tabcolsep}{2.2pt}
\begin{tabular}{lccc}
\toprule
\textbf{Clause (Tier)} & \textbf{Threshold} & \textbf{C2-Original} & \textbf{C2-Replay (Patched)} \\
\midrule
Instruction (1) & $\ge 0.9500$ & 0.9862 (Pass) & 0.9838 (Pass) \\
Structure (1) & $\ge 0.5000$ & 0.5200 (Pass) & 0.5170 (Pass) \\
Retention $R$ (2) & $\ge 0.9800$ & \textbf{0.9666 (Fail)} & \textbf{1.0214 (Pass)$^*$} \\
$\Delta$ Hallucination (2) & $\le +0.0150$ & +0.0102 (Pass) & \textbf{$-$0.0178 (Pass)$^{**}$} \\
Latency Obs. (3) & $\le 3.0000$\,s & 2.4500 (Pass) & 2.4500 (Pass) \\
Peak VRAM (3) & $\le 4.0000$\,GB & 3.8000 (Pass) & 3.8000 (Pass) \\
\midrule
\textbf{Gate Decision} & All Pass & \textbf{Blocked (Open)} & \textbf{Promoted (Closed)} \\
\bottomrule
\multicolumn{4}{l}{\scriptsize $^*p = 0.0182$ (Wilcoxon); $^{**}p = 0.0084$ (Paired $t$-test vs. Baseline).}
\end{tabular}
\end{table}

\subsection{RQ1: Does the gate detect backward regression?}
Yes. When candidate C2-Original was evaluated on the canary, its overall score fell from 0.5693 to 0.5554 ($R = 0.9666$). Violating the Tier 2 threshold ($R \ge 0.9800$), the circuit breaker tripped to OPEN. Traffic remained locked on C1 (SHA-256: \texttt{0ac8c1c0...}), shielding users from regression.

\subsection{RQ2: Is the oracle fast and reproducible?}
Yes. Pre-flight verification executed in $47.02$\,ms on the evaluation harness. Single-sample inference latency averaged $2.45$\,s and peak GPU memory was $3.80$\,GB, satisfying Tier 3 serving SLAs ($\le 3.0$\,s, $\le 4.0$\,GB). Paired testing confirms that observed retention degradations in lexical metrics (ROUGE-1 $p = 0.0372$, BLEU $p = 0.0231$) are statistically significant, ensuring the oracle acts on genuine distributional degradation rather than random measurement noise.

\subsection{RQ3: Does replay restore the gate conditions?}
Yes. Remediation via 10\% selective replay produced C2-Replay (SHA-256: \texttt{28b59045...}). As shown in Tables~\ref{tab:metrics} and \ref{tab:results}, C2-Replay cleared all tiers: overall score rose to 0.5869 ($R = 1.0214 \ge 0.9800$, $p = 0.0182^*$), while hallucination dropped to 0.1717 ($\Delta H = -0.0178$, $p = 0.0084^{**}$). The controller executed an atomic manifest swap and returned to CLOSED, safely promoting the patch with zero downtime.

\section{Threats to Validity}
\textit{Construct validity.} We operationalize behavioral non-regression through nine deterministic metrics rather than human evaluators. While heuristic metrics cannot capture all semantic subtleties, they provide repeatable, zero-cost, millisecond-scale gates suitable for continuous deployment. Pinned canary slices isolate evaluation sets from fine-tuning streams, preventing data leakage.

\textit{Internal validity.} Confounding factors include canary sample composition and replay seed variance. We mitigated selection bias by pinning evaluation IDs prior to fine-tuning and running paired statistical tests across all 3,500 streaming records. The observed drops and recoveries remained consistent across parametric and non-parametric tests ($p < 0.05$).

\textit{External validity.} Experiments were conducted on an event-driven distillation pipeline pairing Alpaca-7B with Gemma-3-1B under customer-service query distributions. While contract specifications and state machine logic are architectural and model-agnostic, validating performance on other LLM families, multi-turn dialogues, and non-English domains remains future work.

\section{Future Plans}
Four research thrusts follow from this prototype:
\begin{enumerate}
\item \textbf{Multi-turn and tool-use contracts:} Add dialogue-state, tool-selection, and argument-schema checks using code and API benchmarks such as HumanEval and ToolBench \cite{DBLP:journals/corr/abs-2107-03374,ICLR2024_28e50ee5}.
\item \textbf{Adaptive replay selection:} Compare uniform 10\% sampling with gradient- or Fisher-informed selection, measuring retention, utility, and compute overhead \cite{NEURIPS2019_fa7cdfad,NIPS2017_f8752278}.
\item \textbf{Distributed consensus:} Coordinate candidate manifests and promotion state across brokers to prevent split-brain releases in multi-node clusters.
\item \textbf{Human triage:} Escalate ambiguous contract failures to human reviewers while retaining automated safety floors.
\end{enumerate}

\section{Related Work}
Continual-learning methods address interference during training: EWC regularizes weight updates, GEM uses episodic memory, and rehearsal revisits prior examples \cite{kirkpatrick2017overcoming,NIPS2017_f8752278,NEURIPS2019_fa7cdfad}; LoRA parameterizes fine-tuning rather than specifying release policies \cite{hu2022yelong}. Our gate is complementary: it inspects candidate artifacts after training and governs promotion.

Production ML testing work provides our engineering foundation: the ML Test Score structures testing rubrics, while systems research highlights technical debt in deployed ML \cite{8258038,NIPS2015_86df7dcf,8804457}. Design by contract motivates formal pre/post-conditions, and circuit breakers isolate failing services \cite{meyer1997object,nygard2007release,10.1145/3542929.3563466}. Holistic benchmarks like HELM and TruthfulQA guide multi-metric evaluation \cite{https://doi.org/10.1111/nyas.15007,lin-etal-2022-truthfulqa,bang-etal-2025-hallulens}; our gate provides operational screening for safe continuous delivery.

\section{Conclusion and Open Science}
Continuous fine-tuning pipelines require behavioral contract gates to protect against silent regressions. We introduced a prototype deployment gate combining a three-tier contract specification, cryptographic SHA-256 weight binding, atomic crash-consistent promotion, and an automated circuit breaker with selective replay remediation. Preliminary evaluation on a 100-example canary demonstrates automated regression interception and recovery.

\textit{Open Science.} A complete, anonymized replication package containing source code, evaluation datasets, metric implementations, audit logs, and statistical analysis scripts is prepared for artifact evaluation to support full community reproducibility \cite{10.1145/3287560.3287596}.

\newpage
% Page 5: References ONLY (ICSE NIER strict 4+1 page rule)
\renewcommand{\IEEEbibitemsep}{0.5pt plus 0.25pt}
\bibliographystyle{IEEEtran}
\bibliography{sample-base}

\end{document}
